%=====================================================================
\chapter{Evaluating Machine Learning Research}\label{ch:mleval}
%=====================================================================
\locator{4}{Part~4 --- Measurement, Data and Analysis}

\begin{outcomes}
By the end of this chapter you will be able to:
\begin{tight}
  \item \textbf{Audit} a pipeline for the eight forms of data leakage.
  \item \textbf{Choose} a validation scheme that matches how the model would be
    deployed.
  \item \textbf{Select} metrics appropriate to the base rate and the costs of
    error.
  \item \textbf{Compare} classifiers across datasets without running twenty
    t-tests.
  \item \textbf{Design} an ablation that shows which component earns the gain.
  \item \textbf{Document} a model and a dataset to a recognised standard.
\end{tight}
\end{outcomes}

\begin{prereq}
\chref{inference} for multiple comparisons and intervals; \chref{measurement}
for construct validity, which applies to labels as much as to questionnaires;
\chref{benchmarking} for measurement under non-determinism.
\end{prereq}

\begin{keyterms}
data leakage $\bullet$ train--validation--test split $\bullet$ temporal split
$\bullet$ group split $\bullet$ baseline $\bullet$ class imbalance $\bullet$
precision and recall $\bullet$ ROC-AUC and PR-AUC $\bullet$ calibration
$\bullet$ decision threshold $\bullet$ ablation $\bullet$ Friedman test
$\bullet$ test set reuse $\bullet$ model card $\bullet$ datasheet
\end{keyterms}

\section{The Claim a Model Makes}

A machine learning paper claims that a model will perform at some level on data
it has not seen, drawn from some population, under some deployment conditions.
Everything in this chapter is about whether the evaluation actually supports
that claim --- and the dominant finding of the last few years is that very
often it does not.

\begin{paperbox}[Leakage across an entire literature]
Sayash Kapoor and Arvind Narayanan, \emph{Leakage and the reproducibility
crisis in machine-learning-based science}, Patterns 4(9),
2023~\cite{kapoor2023leakage}.\oa{}

\medskip
\textbf{Why it matters.} They surveyed 17 fields using ML for scientific
claims, found leakage affecting 294 papers across them, and showed that
correcting it frequently eliminated the reported advantage of machine learning
over much simpler methods. This is not a warning about a rare bug; it is
evidence of a systematic failure mode.

\medskip
\textbf{What to extract.} Their taxonomy of eight leakage types, and their
model info sheet. Fill the sheet in for your own work --- it takes an hour and
is the cheapest insurance available to a computational thesis.
\end{paperbox}

\section{Eight Ways Information Leaks}

\begin{definitionbox}[Data leakage]
\term{Leakage} occurs when information that would not be available at
prediction time is used during training or model selection. The result is
performance that is real on your test set and unattainable in deployment.
\end{definitionbox}

\begin{center}
\small
\begin{longtable}{@{}L{3.4cm}L{4.6cm}L{6.2cm}@{}}
\toprule
\textbf{Type} & \textbf{What happens} & \textbf{Fix} \\
\midrule
\endfirsthead
\toprule
\textbf{Type} & \textbf{What happens} & \textbf{Fix} \\
\midrule
\endhead
\bottomrule
\endfoot
No test set at all
  & Performance reported on training data
  & Hold out a test set and touch it once \\
Preprocessing on all data
  & Scaling, imputation or encoding fitted before splitting
  & Fit every transform inside the training fold. Use a pipeline object so this
    cannot be forgotten \\
Feature selection on all data
  & Features chosen using the full dataset's labels
  & Select inside the cross-validation loop. This one routinely produces
    dramatic apparent performance from pure noise \\
Duplicates
  & Near-identical records in train and test
  & Deduplicate before splitting, including near-duplicates \\
Temporal leakage
  & Training on data from after the test period
  & Split by time, always, whenever the deployment is a forecast \\
Group leakage
  & The same patient, developer or repository in both splits
  & Split by group, not by row \\
Sampling bias in the test set
  & The test set is not drawn from the deployment population
  & Construct it to match deployment, and say how \\
Tuning on the test set
  & Hyperparameters or thresholds chosen by test performance
  & Three-way split, or nested cross-validation \\
\end{longtable}
\end{center}

\begin{worked}[A leakage audit that changes the conclusion]
\textbf{The study.} Predicting defect-prone files from code metrics across 30
repositories. Reported: 0.91 AUC with a gradient boosting model, ``substantially
outperforming logistic regression''.

\medskip
\textbf{Audit finding 1 --- group leakage.} Rows were split randomly. Files from
the same repository appear in both training and test. Repositories differ
enormously in defect base rate, so the model can identify the repository from
its metric profile and predict that repository's base rate. It is not
predicting defects; it is memorising repositories.

\emph{Fix:} split by repository. Train on 24, test on 6, and rotate.

\medskip
\textbf{Audit finding 2 --- temporal leakage.} Files were labelled defective if
a defect was ever found in them, including after the metrics snapshot. The
model uses present metrics to predict labels partly determined by the future,
which is unavailable at deployment.

\emph{Fix:} fix a snapshot date; label from defects found in a defined window
after it.

\medskip
\textbf{Audit finding 3 --- preprocessing leakage.} Features were standardised
across the full dataset before splitting, so test-set statistics informed the
training scaling. Small effect here, but free to fix.

\medskip
\textbf{Audit finding 4 --- no trivial baseline.} With a 12\% defect rate,
``always predict not defective'' achieves 88\% accuracy. The paper reports AUC,
which is better, but never compares against predicting from file size alone ---
which in defect prediction is a notoriously strong baseline.

\medskip
\textbf{After correction.} Grouped, temporal split, pipeline-encapsulated
preprocessing, and a size-only baseline included. Illustrative outcome: AUC
falls to about 0.68, the size-only baseline reaches 0.65, and logistic
regression on the same features reaches 0.67. The honest conclusion is that
gradient boosting offers little over a simple model, and that most of the
signal is size.

\medskip
That is a publishable finding --- it is a negative result about a widely
assumed advantage --- and it is a far better thesis than the original.
\end{worked}

\section{Validation Must Match Deployment}

\begin{center}
\small
\begin{tabular}{@{}L{3.4cm}L{5.0cm}L{5.9cm}@{}}
\toprule
\textbf{Deployment} & \textbf{Correct scheme} & \textbf{What random $k$-fold gets wrong} \\
\midrule
Predict for new instances from the same population
  & Random $k$-fold
  & Nothing. This is the case it is designed for \\
Predict the future
  & Time-series split; train on past, test on future, rolling forward
  & Trains on the future, which inflates results and cannot be reproduced in
    deployment \\
Predict for a new group --- new project, hospital, user
  & Grouped $k$-fold, or leave-one-group-out
  & Lets the model memorise group identity \\
Predict under distribution shift
  & Test explicitly on a shifted set
  & Reports performance for a world that will not occur \\
\bottomrule
\end{tabular}
\end{center}

\begin{alertbox}[The question that settles the scheme]
Ask: \emph{at prediction time in the real system, what would the model know?}

\medskip
Then construct the validation so it knows exactly that and no more. Every
leakage type above is a case of the model knowing something the deployed system
would not. This single question catches most of them without needing to
memorise the taxonomy.
\end{alertbox}

\section{Metrics, and the Base Rate}

\begin{center}
\small
\begin{longtable}{@{}L{3.0cm}L{4.6cm}L{6.6cm}@{}}
\toprule
\textbf{Metric} & \textbf{Answers} & \textbf{Fails when} \\
\midrule
\endfirsthead
\toprule
\textbf{Metric} & \textbf{Answers} & \textbf{Fails when} \\
\midrule
\endhead
\bottomrule
\endfoot
Accuracy
  & What fraction is correct?
  & Classes are imbalanced. At a 1\% base rate, 99\% accuracy is achieved by
    predicting nothing \\
Precision
  & Of those flagged, how many are real?
  & Used alone --- it says nothing about what was missed \\
Recall
  & Of the real ones, how many were found?
  & Used alone --- flagging everything gives recall 1.0 \\
$F_1$
  & Harmonic mean of the two
  & Precision and recall have different costs, which they usually do.
    Weight them by cost instead \\
ROC-AUC
  & Ranking quality across thresholds
  & Heavy imbalance --- ROC looks optimistic because true negatives are
    plentiful~\cite{saito2015precision} \\
PR-AUC
  & Ranking quality, precision against recall
  & Preferred under imbalance. Note the baseline is the positive rate, not 0.5 \\
Calibration
  & Are the probabilities meaningful?
  & Ignored almost always, and essential if a probability feeds a decision \\
\end{longtable}
\end{center}

\begin{conceptbox}[The threshold is a decision, not a default]
Classifiers output scores; 0.5 is a convention with no claim to optimality. The
right threshold follows from the costs of the two errors.

\medskip
In defect prediction, a false negative means a defect ships and a false
positive means a developer spends twenty minutes checking clean code. Those are
not equal, and the ratio between them determines the threshold. State the cost
ratio you assume, choose the threshold from it on \emph{validation} data, and
report performance at that threshold on test data.

\medskip
A paper reporting only threshold-free metrics has evaded the question its users
must answer.
\end{conceptbox}

\section{Baselines and Ablations}

\begin{center}
\small
\begin{tabular}{@{}L{3.4cm}L{10.9cm}@{}}
\toprule
Trivial baseline    & Majority class; random; predict the mean. Establishes the
                      floor, and occasionally beats the proposed method \\
Simple baseline     & Logistic regression, a decision tree, or one strong
                      feature. In software engineering, file size is a famously
                      strong baseline for defect prediction \\
Current best        & The actual state of the art, tuned with the same budget
                      you gave your own method. Tuning yours for a week and the
                      baseline for an hour is not a comparison \\
Ablation            & Your method with one component removed at a time. This is
                      what identifies which part earns the gain \\
\bottomrule
\end{tabular}
\end{center}

\begin{pitfall}[The unablated improvement]
A paper proposes a method combining a new architecture, a new loss, a new
augmentation scheme and a longer training schedule, and reports a 2-point gain.
Without ablations, no reader --- including the authors --- knows which of the
four produced it, or whether it was the longer schedule alone.

\medskip
This matters beyond credit assignment: an unablated improvement cannot be
transferred, because nobody knows what to transfer. Run the ablations, and
report the ones where removing a component changed nothing. That is
information, and omitting it is what makes a paper unreproducible in the useful
sense.
\end{pitfall}

\section{Comparing Across Datasets}

Comparing $k$ classifiers over $n$ datasets with pairwise t-tests is wrong on
two counts: the multiplicity is uncorrected (\chref{inference}), and the
normality assumption is untenable for performance scores across heterogeneous
datasets.

\begin{conceptbox}[The standard procedure]
Dem\v{s}ar's recommended approach~\cite{demsar2006statistical} is now
conventional and is two steps:

\begin{tightnum}
  \item \textbf{Friedman test} on the ranks of the classifiers across datasets.
    Non-parametric, and it asks the right global question: are these
    classifiers performing equivalently overall?
  \item If significant, a \textbf{post-hoc test} --- Nemenyi for all-pairs
    comparison, or Bonferroni--Dunn when comparing everything against one
    control method, which is usually what you actually want.
\end{tightnum}

\medskip
Report the result as a \textbf{critical difference diagram}: classifiers on a
rank axis, with bars joining those not significantly different. It communicates
the whole comparison in one figure and makes over-claiming difficult.
\end{conceptbox}

\begin{lstlisting}[language=Rlang, caption={Comparing four classifiers across twelve
datasets, the way it should be done.}]
# rows = datasets, columns = classifiers, cells = performance
friedman.test(as.matrix(results))

# If significant, post-hoc pairwise comparison on ranks
library(PMCMRplus)
frdAllPairsNemenyiTest(as.matrix(results))

# Report mean ranks alongside, not raw means: ranks are what was tested
colMeans(t(apply(-results, 1, rank)))
\end{lstlisting}

\section{Test Set Reuse}

Every time you look at test performance and change something, you leak a little
information. Over a project --- or over a community's decade of using one
benchmark --- this accumulates into overfitting the test set, and reported gains
stop transferring.

\begin{center}
\small
\begin{tabular}{@{}L{3.6cm}L{10.7cm}@{}}
\toprule
Three-way split      & Train, validation, test. All decisions on validation;
                       test touched once, at the end \\
Nested cross-validation & When data is scarce: an inner loop for tuning, an
                       outer loop for estimating performance \\
Count the looks       & Record how many times you evaluated on test. If it is
                       more than once, say so --- the number is meaningful \\
A fresh test set      & The only real remedy for a community-level problem.
                       Where a new one exists, report on it \\
\bottomrule
\end{tabular}
\end{center}

\section{Documenting Models and Data}

\begin{center}
\small
\begin{tabular}{@{}L{3.2cm}L{11.1cm}@{}}
\textbf{Datasheet}~\cite{gebru2021datasheets}
  & For a dataset: motivation, composition, collection process, preprocessing,
    uses, distribution, maintenance. Crucially, its known biases and the uses
    for which it is \emph{not} appropriate \\
\textbf{Model card}~\cite{mitchell2019model}
  & For a model: intended use, out-of-scope use, training data, evaluation
    data, metrics \emph{disaggregated by relevant subgroup}, ethical
    considerations, caveats \\
\end{tabular}
\end{center}

Disaggregated evaluation is the part most often skipped and most often
consequential: an aggregate accuracy of 0.90 can conceal 0.95 on one subgroup
and 0.60 on another, and only reporting by subgroup reveals it.

\begin{traditions}{evaluating a model}
\tradEMP{A model evaluation is an experiment. It needs repetitions across
seeds, intervals rather than point estimates, a baseline, and correction for
the comparisons made.}
\tradDSR{The model is the artefact. Held-out performance is artificial
summative evaluation; whether anyone benefits from deploying it is a separate
and usually unasked question (\chref{designscience}).}
\tradQUAL{Asks what the labels mean and who produced them. A defect label is a
human judgement, and label noise is a measurement validity problem
(\chref{measurement}), not merely a data quality nuisance.}
\tradFORM{Provides generalisation bounds and guarantees under stated
assumptions --- typically far looser than observed performance, but true for
reasons that do not depend on a test set.}
\end{traditions}

\begin{lab}[Audit a pipeline]
Take any ML pipeline --- yours, a published replication package, or a public
notebook.

\begin{tightnum}
  \item Work the eight leakage types. For each, state where in the code it
    would occur and whether it does.
  \item Ask the deployment question and check the validation scheme matches.
  \item Add a trivial baseline and one strong simple baseline. Record the gap.
  \item Re-run with the leakage fixed. Record how much performance falls.
  \item Write the model card, including disaggregated metrics for at least one
    meaningful subgroup.
\end{tightnum}

\medskip
Step 4 is the one to report in the seminar. A drop is not a failure; it is the
study working.
\end{lab}

\begin{checkpoint}
\begin{tightnum}
  \item Why does fitting a scaler before splitting constitute leakage, and how
    large is the problem in practice?
  \item You are predicting whether a developer will abandon a project. What
    split do you use, and why is random $k$-fold wrong?
  \item At a 2\% positive rate, why is ROC-AUC misleading, and what should you
    report?
  \item Why are pairwise t-tests wrong for comparing five classifiers across
    fifteen datasets, and what replaces them?
\end{tightnum}
\end{checkpoint}

\begin{chaptersummary}
\begin{tight}
  \item Leakage is systematic, not rare, and correcting it frequently
    eliminates reported advantages of complex methods over simple ones.
  \item Ask what the deployed system would know at prediction time, and build
    validation to match. That question catches most leakage.
  \item Accuracy fails under imbalance; ROC-AUC flatters it; PR-AUC and
    calibration are usually what you want. The threshold follows from costs.
  \item Include a trivial baseline, a strong simple baseline, and a fairly
    tuned state of the art. Ablate, and report the components that changed
    nothing.
  \item Compare across datasets with Friedman plus a post-hoc test, reported as
    a critical difference diagram.
  \item Touch the test set once, and count the times you did not.
  \item Publish a datasheet and a model card, with metrics disaggregated by
    subgroup.
\end{tight}
\end{chaptersummary}

\begin{reviewq}
  \item Kapoor and Narayanan found leakage across 17 fields. What features of
    how ML is taught and published produce this, and which would be easiest to
    change?
  \item Benchmark test sets are reused by whole communities for years. Propose a
    workable governance mechanism for a benchmark in your area, and identify
    who bears its cost.
  \item Defect labels come from issue trackers and are noisy and biased
    (\chref{quasiexp}). What does label noise do to a reported AUC, and how
    would you estimate its magnitude?
  \item \reg{PhD Regs 2024}{\S14.3(B)} asks whether results are replicable. For
    a deep learning dissertation, enumerate exactly what would have to be
    deposited for a third party to replicate the reported numbers, and estimate
    what fraction of published work meets that bar.
\end{reviewq}
